\begin{table}[t]
\centering
\caption{Standardized Effect Sizes}
\label{tab:sde}
\begin{tabular}{lcccccc}
\toprule
Outcome & $\hat{\beta}$ & SE & SD($Y$) & SDE & SE(SDE) & Classification \\
\midrule
TA per 1,000 & 0.312 & 0.177 & 2.200 & 0.1418 & 0.0804 & Moderate positive \\
Log(TA + 0.01) & 0.0034 & 0.0521 & 1.553 & 0.0022 & 0.0336 & Null \\
\bottomrule
\end{tabular}
\begin{tablenotes}[flushleft]\footnotesize
\item \textit{Notes:} \textbf{Country:} United Kingdom. \textbf{Research question:} Does reducing the household benefit cap increase local authority temporary accommodation burdens? \textbf{Policy mechanism:} The November 2016 benefit cap reduction from \pounds26,000 to \pounds20,000/\pounds23,000 limited total welfare payments for out-of-work households, potentially forcing rent shortfalls and housing displacement. \textbf{Outcome definition:} Households in temporary accommodation per 1,000 population at end of financial year, from MHCLG statutory homelessness Table 784. \textbf{Treatment:} Continuous; capped households per 1,000 population at May 2017. \textbf{Data:} MHCLG Table 784 (2013--2018), DWP benefit cap statistics, NOMIS population estimates; 278 English local authorities observed annually. \textbf{Method:} TWFE with LA and year fixed effects; standard errors clustered at LA level. \textbf{Sample:} English local authorities with non-missing TA data across all six years; 1,592 LA-year observations. SDE $= \hat{\beta} / \text{SD}(Y)$ where SD($Y$) is the pre-treatment standard deviation. Classification refers to magnitude, not statistical significance: Large ($|$SDE$| > 0.15$), Moderate ($0.05$--$0.15$), Small ($0.005$--$0.05$), Null ($< 0.005$).
\end{tablenotes}
\end{table}
